\section{Prices and local derivatives of the marked activities}
\label{sec:marked-activity-jets}
The marked integral permits a bounded measurable insertion. Its
coarse local extension nevertheless has the derivatives needed for
the projection argument. The reason is the compensated transport:
it holds the fine field read by the insertion fixed. Only its smooth
cutoff and the surrounding integral are differentiated.

\subsection{Paying for the distinguished letter}
Let $X$ be the fine insertion support, $c=C_{B_f}(X)$ its fine
cover number, $s$ the fine support exponent, and $w$ its uniform
bound. Its fine mass is $m=e^{sc}w$. Fix the support reserve
$\varepsilon>0$, with $s>4\varepsilon$, and let $g>0$ be the
normalization gain. Introduce
\begin{equation}\label{eq:marked-debt-price}
 D=\frac{m e^{-\varepsilon c}}g,\qquad
 J=g e^{-(s-\varepsilon)c},\qquad DJ=w.
\end{equation}
The last equality follows by cancellation of the two exponentials.
It also covers $w=0$. We give the marked letter price $J$ in the
mixed component estimate and keep the factor $D$ outside it.

Let $X'$ be its full enlarged coarse support. The footprint
estimate of Section~\ref{sec:inherited-activity-prices} has the form
\[
 C_B(X')\le C_0+\frac{10C_0}{L}c,
\]
where $C_0$ depends only on the fixed enlargement depth. Choose
the block ratio so that $10uC_0/L\le\varepsilon$, where $u$ is the
mixed estimate's support exponent. The marked contribution to
the elementary price, including the two length reserves, then obeys
\begin{align*}
 J e^{2\varepsilon c}e^{uC_B(X')}
 &\le g e^{uC_0}e^{-(s-4\varepsilon)c}\\
 &\le g e^{uC_0}.
\end{align*}
At any root there is at most this one new letter. Thus replacing
the old price at that address by $J$ increases both the optional
root bound and the total root bound by at most $g e^{uC_0}$.
All other prices and all supports remain unchanged. Applying the
mixed connected-component estimate of
Section~\ref{sec:mixed-activity-prices} to this enlarged price
family bounds the marked components, including any excess-cardinality,
support-size and forced-count restrictions required there.

This enlargement is purchased at the same scale as the background.
If its original root price and $g$ have negative upper exponent
$q<0$, their sum has that exponent too. The geometric degree factor
has upper exponent zero. It is therefore eventually at most
$\beta_f^{q/2}$; this quantity tends to zero and may be used as
the new tree ratio. This follows directly from the strict exponent
comparison in Section~\ref{sec:leaf-expansion}. It makes no new
choice of a history or a fine-to-coarse transition.

\subsection{Holding a merely measurable insertion fixed}
Fix a component $P$ containing the marked letter, a healthy coarse
configuration on its full support, and a smooth coarse probe based
there. Use the compensated transport from
Section~\ref{sec:compensated-variation}. For fixed old Haar variables
and Gaussian point $z$, let $U_0(z)$ be the physical fine field at
the base point. On the nonzero support of the transported marked
cutoff, the transport identity gives agreement with $U_0(z)$ on
the inspected support of $h$. Since $h$ is local, the marked
factor becomes
\begin{equation}\label{eq:marked-held-factor}
                  h(U_0(z))\,\zeta_\star(y,z).
\end{equation}
If the cutoff is zero, both expressions are zero. The identity
holds as a germ in the coarse probe parameter $y$. The truncation
in the compensated transport preserves this germ, because its
threshold is larger than the cutoff support. This is exactly the
cutoff comparison used for the unmarked transported factors.

The function $h(U_0(z))$ is constant in $y$ and has absolute value
at most $w$. The cutoff derivative estimate therefore gives
amplitude $w$ and derivative scale at most $J_0\sigma_\star$, where
$\sigma_\star$ is the marked letter's support scale and $J_0$
has upper exponent zero. No derivative of $h$ has been taken.
This observation also gives the locally uniform integrable
derivative envelopes: multiply the existing smooth cutoff envelopes
by the bounded fixed function $h(U_0(z))$.

Write $a_i$ for the unmarked elementary jet amplitudes and $r$
for their common scale factor. The sum of the support scales in
the component is at most twice its geometric size $W_P$.
Leibniz's rule applied to \eqref{eq:marked-held-factor} and the
unmarked factors gives amplitude
\[
                    A_P=w\prod_{i\in P\setminus\{\star\}}a_i
\]
and derivative scale at most $2(r+J_0)W_P$.
The same joint warning event can be retained. Outside that event,
one warning factor has the zero germ supplied by the strict
small-fluctuation plateau. The marked letter is a regular or forced
insertion letter, never a warning letter; erasing it leaves the
warning set unchanged. Thus the zero germ still occurs among the
remaining factors.

Now include the Jacobian and projected Gaussian density of
Section~\ref{sec:gaussian-differentiation}. Their joint derivative
scale is bounded by $U_0\mathcal F_I(z)^k$, with amplitude one,
where $\mathcal F_I\ge1$ is the Gaussian coordinate polynomial
on the actual local read set $I$. For derivatives through order
$j$, the full integrand has the bound
\begin{equation}\label{eq:marked-pointwise-jets}
 \mathbf1_{\mathcal W_P}(z) A_P
       [1+U_0+2(r+J_0)]^j(1+W_P)^j\mathcal F_I(z)^{kj}.
\end{equation}
To verify the scale conversion, use
\[
 1+U_0\mathcal F_I^k+2(r+J_0)W_P
 \le[1+U_0+2(r+J_0)](1+W_P)\mathcal F_I^k
\]
and raise both nonnegative sides to the integer power $j$.
The estimates concern the unmasked extension on a neighborhood
of the healthy base point; they do not differentiate the flag mask.

\subsection{Integrating the derivative bound}
The Gaussian moment and warning estimate of
Section~\ref{sec:gaussian-moment-prices} gives, for some fixed
integer $K$ and scalar $F$ of upper exponent zero,
\[
 \int\mathbf1_{\mathcal W_P}\mathcal F_I^{kj}\,d\gamma
          \le (1+|F|)(1+W_P)^K\rho^{N_P}.
\]
Here $N_P$ is the number of warning letters and $\rho$ is the
half-warning factor of that estimate. Explicitly, Cauchy--Schwarz
uses the moment of degree $2kj$ and the square root of the joint
warning probability. This explains both the doubled moment degree
and the preserved half-warning gain. The old Haar measure has total
mass one, so introduces no further volume factor.

Integrating \eqref{eq:marked-pointwise-jets} therefore bounds the
local derivative amplitude by
\begin{equation}\label{eq:marked-integral-jets}
 B_* A_P\rho^{N_P}(1+W_P)^{K+j},\qquad
 B_*=(1+|F|)[1+U_0+2(r+J_0)]^j.
\end{equation}
The finite-order differentiation-under-the-integral argument of
Section~\ref{sec:gaussian-differentiation} applies: the relative
Gaussian density supplies an integrable local envelope, the
transported insertion is bounded and constant in the probe, and
all remaining factors have the already proved smooth envelopes.
Thus this is a bound for derivatives of the actual local integral.
The constants $j,k,K$ are fixed before the coupling limit. Products,
sums and fixed powers show that $B_*$ has upper exponent zero.

Let $v_P$ be the mixed component price with the marked elementary
price replaced by $J$, with the warning gain and the size weight
$(1+W_P)^{K+j}$ included. The exact identity $DJ=w$ from
\eqref{eq:marked-debt-price} rewrites
\eqref{eq:marked-integral-jets} as $D B_*v_P$.
This proves the marked analogue of the unmarked local jet estimate
from the same sector data, rather than postulating a new regularity
property of $h$.

To make a globally measurable gauge-invariant extension, first
average the local projected integral over the original finite tag
choices, then average its gauge translates against product
probability Haar measure. Locality is preserved because only the
gauge coordinates incident to its read support matter. Haar
invariance makes the result gauge invariant. At a healthy base
point, the local covariance and chart bounds persist on a probe
neighborhood. Change of variables in the local integral shows
that it is already gauge invariant there. Hence the gauge average
agrees with that local integral as a germ, and the derivative bound
is unchanged. Independent forest averaging also preserves the bound.
We have constructed a measurable local extension $g_P$ satisfying
\[
 b_P(V)=\mathbf1_{\{X_P\cap\operatorname{Flags}(V)=\varnothing\}}
                                  g_P(V)
\]
for regular components, with derivatives through the purchased
order bounded by $D B_*v_P$ at healthy configurations. This is
finite-order smoothness on the stated domain.

\subsection{The logarithmic derivative and its cluster prices}
Let $f_P$ be the unmarked local extension. For a regular cluster
tuple $\Gamma=(P_1,\ldots,P_n)$ its unmasked marked logarithmic
term is
\begin{equation}\label{eq:marked-local-log-term}
 h_\Gamma=\frac{u(P_1,\ldots,P_n)}{n!}
       \sum_{r=1}^n g_{P_r}\prod_{s\ne r}f_{P_s}.
\end{equation}
This is the derivative at zero of the corresponding product of
affine tilted factors. All terms have the mask of the union
support. Indeed on a healthy union every constituent mask is one;
if the union contains a flag, each summand has a zero factor,
either its marked factor or one of its ordinary factors.

Suppose $f_P$ has local amplitude $A_P^0$ and $g_P$ has amplitude
$D A_P^1$, both at derivative scale one. Put
$A_P=A_P^0+2A_P^1$ and inflate this price to
$\widehat A_P=2^{j+2}A_P$. On a healthy union each summand in
\eqref{eq:marked-local-log-term} has amplitude at most
$D\prod_r A_{P_r}$ and derivative scale $n$. Its derivatives
through order $j$ are consequently bounded by
\[
 D\frac{|u(P_1,\ldots,P_n)|}{n!}
                     n(1+n)^j\prod_r A_{P_r}
 \le D\frac{|u(P_1,\ldots,P_n)|}{n!}
                                      \prod_r\widehat A_{P_r}.
\]
For the last inequality use $1+n\le2^n$ and $n\le2^{2n}$.
Thus the canonical cluster majorant at the inflated prices pays
all these derivatives as well as the scalar tilt derivative.
The rooted cluster estimate and its tail then give
\eqref{eq:leaf-cluster-tail}. The marked mixed-component estimate
proved above supplies its order-zero counterpart, with the
excess-cardinality restriction for nonprincipal terms. The
polynomial support-weight conversion and strict exponent reserve
in Section~\ref{sec:leaf-expansion} therefore apply to these
actual local terms.

For forced components use their absolute marked integral price
from \eqref{eq:marked-integrated-price} and the same mixed estimate.
The affine joint bounds propagate through the exclusion
regrouping by the product rule, keeping the exponential
multiplicity factorials. Section~\ref{sec:forced-exclusion}'s
rooted argument applies to these inflated prices and gives the
summable joint forced-price family required in
\eqref{eq:leaf-season-bounds}. Every nonzero marked component
contains the projected insertion anchor, so the derivative vanishes
at all other supports. This proves the structural and local
analytic inputs to the marked differentiation. Their simultaneous
scale purchases are made as follows.

\subsection{One coupling threshold for the marked estimates}
Fix the derivative and spatial orders, the target support exponent
$a\ge0$, and a reserve $e>0$ with $a+2e$ no larger than the
available mixed support exponent. Let $\mathcal O_n$ be the
unmarked rooted component price with size weight $n$, and $C_n$
the corresponding marked mixed price after adjoining $J$. Both
have upper exponent $q<0$. Write $C_n^{(1)}$ for the marked price
with excess cardinality at least one; its extra tree factor gives
upper exponent $3q/2$. Put
\[
 t=(1+2^{j+2}B_*)(\mathcal O_n+2C_n),\qquad
 v=\mathcal O_n+2C_0,\qquad u=10(t+v),
 \qquad \gamma=\beta_f^{q/2}.
\]
These quantities are positive in the selected construction. The
inflated regular prices above have rooted allowance at most $t$.
The forced prices have allowance at most $v e^{-M}$, where $M$
is the compulsory forced mark. These assertions follow by adding
the unmarked and marked mixed bounds, then multiplying the
regular bound by $2^{j+2}B_*\le1+2^{j+2}B_*$.

The functions $t,v,u$ have upper exponent $q$. Since the geometric
degree factor $4(2B+1)^4/e$ has upper exponent zero, eventually
\[
 \frac{4(2B+1)^4t}{e}\le\gamma,\qquad
 \frac{4(2B+1)^4u}{e}\le\tfrac12,\qquad
       \gamma\le\tfrac12,\qquad u\le\tfrac12.
\]
At the same scale the shifted fine radius is at most twice the
coarse radius, by the radius estimate of
Section~\ref{sec:ex-lattice-price-receipts}. There are finitely many
eventual inequalities at fixed orders; choose their common
threshold before choosing the finite torus or insertion. The
rooted bounds are uniform in both, so this threshold works for
every permitted insertion on every such torus.

For the second expansion the two rooted contributions, after
inflating copies and extracting the compulsory mark, are bounded by
\[
 4e^M(ve^{-M})+4t(1-\gamma)^{-1}
                   \le4v+8t\le10(t+v)=u.
\]
The connected multiplicity argument in
Section~\ref{sec:forced-exclusion} thus gives forced joint price
at most $2u e^{-M}$. Its differentiated price is $D$ times this
quantity, as required in Section~\ref{sec:leaf-expansion}.

Finally, let $A_{e,r}$ be the polynomial support-weight price in
that section. It has upper exponent zero at fixed orders. The
nonprincipal regular tail divided by the fine mass reserve is
bounded by
\[
 \frac{A_{e,r}}g\left(B_* C_n^{(1)}+
                          t\frac{\gamma}{1-\gamma}\right).
\]
Here $g=\beta_f^q$, so $g^{-1}$ has upper exponent $-q$.
Both numerator terms have upper exponent $3q/2$, since
$(1-\gamma)^{-1}\le2$. The quotient therefore has upper exponent
$q/2$ and is eventually at most $\beta_f^{q/4}$.
This purchases the tail and joint forced prices used by the
actual marked expansion. The final principal-term and projection
allowances must still be compared with the selected outgoing
budget, using the same already chosen scale.
